%% APPENDIX A -- spectral response, frequency conventions, primary-beam response.
%% Owner: appendix-triage.  Carries \label{app:conventions} (cited by
%% 04_method.tex) and \label{app:pbresp}.
\section{Spectral response, frequency conventions and the primary-beam
response}
\label{app:conventions}

This appendix sets out the three things every threshold quoted in this paper
is checked from: the frame of each frequency, the instrumental response, and
the primary-beam factor that turns a measured noise into a flux at the star.

%% The single-column figure of the three-point Hanning response against
%% sub-channel phase is removed; its numbers are in the peak-channel
%% paragraph below.  QUEUE_integrate.md entry S6: make_all.sh must stop
%% running make_fig_hanning.py and figorphan.py must carry
%% hanning_response.pdf in PENDING, or the figure gate fails.

\subsection{Frequency frames, linewidth and smearing}
%% (subsection label retired: nothing cites it)

Every frequency axis here -- searched spectra, tabulated ranges and
crossing frequencies alike -- is \emph{topocentric} as delivered by the
ALMA correlator, and the search applies no velocity-frame shift.
Velocities enter at the vetting stage, where the line mask carries every
crossing frequency to the barycentre at its own block's epoch and then to the
stellar frame before comparing it with a transition
(Appendix~\ref{app:linemask}). The residual error is then the catalogue
radial-velocity uncertainty alone, a few km\,s$^{-1}$ within 40\,pc;
second-order relativistic terms at the \MaskHalfKms\,km\,s$^{-1}$
scale are $\lesssim5$\,kHz at 350\,GHz, below the finest native channel
searched. The stacking is drift-aligned and inverse-variance weighted:
coherent in frequency trajectory, not in electric-field phase.

The native channel width is also the linewidth convention for every EIRP
threshold. The statistic is the single-channel flux density after drift
correction, so EIRP$_{\rm trig}$ is the minimum \emph{total} power of a
channel-confined carrier, flat in the assumed linewidth below the channel and
never a spectral power density; flatness is exact only for a centred tone,
and a boundary-straddling one divides its power between two channels, which
is the loss the injections carry and one reason EIRP values are quoted to two
significant figures. The Hz-resolution
surveys' limits are total powers in the same sense, each the best case of
its own channelisation, which makes the axes of Fig.~\ref{fig:context}(a)
commensurable; and no reprocessing recovers resolution the correlator never
recorded. Channel widths come from the measurement sets' own CHAN\_WIDTH;
MS~v2 defines no ``EFFECTIVE\_BANDWIDTH'' column, so the quoted
channelisations are separations and are not claimed to be noise bandwidths.

ALMA applies online Hanning smoothing by default \citep{ALMAHandbook}, so
the spectral response is ${\sim}2$ channels wide and neighbouring channels
are correlated. For the $0.25/0.5/0.25$ kernel the coefficients
are $\rho_1=\HanRhoOneFrac=\HanRhoOne$ and $\rho_2=\HanRhoTwo$, and the
smoothing narrows the noise by
$\sigma_{\rm smoothed}/\sigma_{\rm raw}=\HanSigRatio$. The correlation is
not confined to adjacent channels: the lag-two term is a sixth of the
lag-one term. The peak
channel holds \HanBest{} of an unresolved carrier's flux at a channel centre
and \HanWorst{} on a boundary, a median \HanMed{} over sub-channel phase.

\emph{That loss enters the limits exactly once, and it is the thresholds
rather than the limits that are free of it.} A window's trigger flux is the
bare $\LtCommonSnr\sigma$ of its own spectrum, so setting the \emph{total}
power of a channel-confined carrier against it would be optimistic by
$\times\HanFacBest$--$\times\HanFacWorst$. The injected carriers are not set
against it in that way: each is deposited with the correlator's own channel
response, at a sub-channel phase drawn afresh, and recovered through the same
trigger, which reads the peak channel. The measured completeness therefore
already contains the loss, and factorises into it exactly:
$\EirpNinetyMultA\,P_{\rm trig}$ is $\times\RpRespFac$ for the channel
response and $\times\RpMarginNinety$ for the margin a peak channel needs
above the nominal trigger for nine carriers in ten to return. So
$P_{\rm trig}$, the sensitivity $S_{\min}$ below and the filled symbols of
Fig.~\ref{fig:context}(a) are nominal throughout, which is what makes them
commensurable with other programmes' detection thresholds, while
$\mathrm{EIRP}_{90}$ carries the correction and is not corrected again.
The deposition falls \InjShortPct{} per cent short of the ideal response
averaged over phase, so the completeness is conservative by that much; and
where in a channel a carrier was placed makes no difference to what is
recovered, because at a typical searched drift rate it crosses about
\InjDriftChan{} channels during one observation and samples the whole
response.

\emph{Intra-integration smearing, which the injections do not model.} Each
injected tone is deposited at one frequency per integration, so the campaign
carries the channel response and the drift between integrations but not the
sweep within one. A carrier moving $\Delta$ channels during an integration
spends its dwell time spread across the delivered channel response, so the
amplitude the correlator accumulates is the mean of that response over the
sweep: $\eta_{\rm smear}$ is that mean, and every limit in this paper is
divided by it. The response is the one that sets $\times\RpRespFac$ above --
Hanning-smoothed, and channel-averaged where the window is one of the
\NmRespNAvgA{} Class~A windows that are. \emph{The sweep cannot act inside
the transform}: these channels are $\NmRespChanTau$ times wider than the
reciprocal of an integration, so a sweep does not destroy coherence but
redistributes dwell time. The transform-limited amplitude
$\mathrm{sinc}(\pi\Delta/2)$, which is what a sweep leaves when the spectral
resolution is set by the integration itself, is the more severe of the two
wherever they differ: in the worst case, a carrier crossing
\NmRespDeltaWorst{} of a channel, it asks for $\times\NmSincFacWorst$ where
the delivered response gives $\times\SmcFacWorst$ -- a correction larger by
\NmRespOverPct{} per cent. The two forms part by more than one per cent in \NmRespNPart{} windows,
and the correction itself exceeds one per cent in \NmRespNApplied.
The $\Delta$ that sets a ninety
per cent point is the one at the ninetieth percentile of each window's own
drift draw, \SmcQninePct{} per cent of the searched ceiling over the
campaign's \SmcNTone{} tones. The integration length is recorded for
\SmcNMeas{} of the \NWindows{} windows, \CsNDumpWin{} of them in the delivered
metadata and \SmNRecov{} more in the exposure records the search kept. For the
other \SmcNBound{} the factor is \emph{bounded} at \SmcDtCeil\,s, the longest
integration delivered anywhere, where none falls below \SmcBoundEta.
Nothing moves in the median window, and \SmcNAbove{} windows on
\SmcNAboveStar{} stars rise by more than one per cent, up to
$\times\SmcFacWorst$: \SmcAboveStars. Re-scoring the \SmcValNWin{} injected
windows' own ladders with every tone de-rated by its own $\eta_{\rm smear}$
gives a penalty never larger than the one applied and smaller by at most
\SmcValWorstPct{} per cent, so the correction errs the safe way.

\subsection{The primary-beam response}
\label{app:pbresp}

The sensitivity quoted for a window is $S_{\min}=5\sigma/A$, where $A$ is the
primary-beam response at the star's offset from the phase centre: the noise
$\sigma$ is measured in apparent flux, the limit is quoted in true flux, and
the two differ by $A$. That offset is not a nuisance. Archival fields are
pointed at a catalogue position, often decades old, while these are
high-proper-motion stars, so the star can sit well off axis. Over the
\PbNWin{} windows the offset has a median of \PbOffMed$''$, a 99th
percentile of \PbOffPNN$''$ and a maximum of \BmOffMax$''$. The response $A$
runs from \BmRespMed{} in the median down to \BmRespMin. It falls below $0.99$
on \PbNLowNN{} windows and below $0.9$ on \BmNBelowNine. Every one
of those is a physically explicable J2000 pointing on a
fast-moving star rather than an error.

\emph{The response applied is a Gaussian of the full width the observatory
documents}, $\RingPbCoef\,\lambda/D$ \citep{ALMAHandbook}, evaluated at each
window's \emph{own} dish diameter. The dish matters: the \PbNAcaRow{} rows
taken on the \PbAcaRatio$\times$ smaller compact-array antennas have a
correspondingly wider beam, and a Gaussian forced to $D=12$\,m reproduces the
per-window geometry only to $\SensPbResidTwelve$ against $\SensPbResid$ for
the per-dish form. That width is a physical choice rather than a convention.
The blocked-aperture response, which carries the shadow of the secondary and
is what \textsc{casa} applies \citep{CASA2022}, has a half-power
width of $\BmAiryFwhm\,\lambda/D$ on a 12\,m antenna and
$\BmAiryFwhmAca\,\lambda/D$ on a 7\,m one. The Gaussian used here is that
beam to \BmDocAiryPct{} per cent and departs from it by at most
\BmModelMaxPct{} per cent anywhere in the sample. The other width available,
a Gaussian of the Airy first-null width, is wider and departs from that beam
by up to \BmModelMaxNullPct{} per cent, which is why the documented width is
the one adopted; what it would do to the limits of the windows furthest off
axis is in \S\ref{sec:pninetysel}.

\emph{A window is used where the two forms agree.} At low response a small
error in the beam model becomes a large error in inferred flux, and the two
forms part company as the offset grows, so the limit on retention is the
offset at which they differ by \BmFloorTolPct{} per cent:
$\BmFloorC\,\lambda/D$ from the pointing centre, where the response is
\BmFloorResp. The furthest window in the primary census lies at
$\BmCensusWorstC\,\lambda/D$, where the two agree to \BmCensusWorstPct{} per
cent, so the criterion excludes nothing from it; it is stated because it is
what admits the one star this survey covers far off axis.

\emph{That star is $\epsilon$\,Eridani, at \BmEpsDistPc\,pc the nearest in the
sample with Band~6 coverage.} \BmEpsNOn{} of its \BmEpsNWin{} windows, in
\BmEpsNBlock{} execution blocks, place it at the pointing centre and need no
correction at all. The other \BmEpsNOff{} lie
$\BmEpsCLo$--$\BmEpsCHi\,\lambda/D$ off it, where the response is
\BmEpsRespLo--\BmEpsRespHi. Those offsets are inside the main lobe -- the
blocked-aperture response does not reach its first null until
$\BmAiryNull\,\lambda/D$ -- and there the two beam forms agree to
\BmEpsAgreePct{} per cent, which is the extra uncertainty those limits carry.
No window reaches the trigger, the largest statistic being \BmEpsTMax. Over
\BmEpsHours\,window-hours the deepest of the \BmEpsNFine{} fine windows
reaches $\BmEpsEirpFine$\,W at
\BmEpsChanFineKHzLo--\BmEpsChanFineKHzHi\,kHz, in \BmEpsExtEb; the deepest
that lies inside the primary census is in \BmEpsCenEb{} at
$\BmEpsCenEirp$\,W, and that is the limit Table~\ref{tab:persys} gives for
this star. The
\BmEpsNCoarse{} coarse ones reach
$\BmEpsEirpCoarseLo$--$\BmEpsEirpCoarseHi$\,W. \BmEpsNCensus{} of these
windows fall inside the primary census and carry its statistics --
\BmEpsNCensusOff{} of them off axis and \BmEpsNCensusOn{} at the pointing
centre -- counted with it in Table~\ref{tab:blockledger}; the other
\BmEpsNExt{} are listed apart from the archival extension's null and are
reported here instead. Nothing about the windows themselves sets them apart:
all \ExNSearched{} were searched with the full control ensemble, the largest
statistic among them is \ExTMax, each carries its own phase centre and
integration time from its measurement set, and the floor above admits every
one of them -- their lowest response is \ExRespMin, and the window furthest
from a pointing centre anywhere in the archival coverage is one of these, at
\ExCWorstExt\,$\lambda/D$.

\emph{Nor does any block of this star carry an annual-parallax record.} The
displacement follows from its catalogued parallax, \BmEpsPlxMas\,mas, which
is the largest in the sample: at the epoch of \BmEpsCenEb{} it is
\ExCenPxArcsec$''$, or \ExCenPxBeams{} of that block's synthesised beam, and
\ExBestPxArcsec$''$, or \ExBestPxBeams{} of a beam, at the epoch of the block
that sets the deepest limit on this star.
That is smaller than the \SensPxSmallBeams{} of a beam at which the retained
amplitude was measured below, so that measurement bounds the loss instead of
being extrapolated to it, and the limit is optimistic by at most
\BmEpsPxBoundPct{} per cent on this account. It is carried as a one-sided
systematic in Table~\ref{tab:persys} rather than folded in, because the
retention every other window's correction uses is measured on that block's
own baselines and these deliveries expose none. The
flux-scale
comparison of \S\ref{sec:method} is a separate continuum measurement on the
same star, a block mean over the \PxEpsN{} compact-array executions that place
it at the pointing centre, \PxEpsHours\,h in Band~\PxEpsBand, where no beam
correction enters either.

\emph{The correction is applied exactly once.} The offset is computed for
every measurement set from its own \texttt{FIELD::PHASE\_DIR} and the
propagated stellar position, so no window carries a zero-by-assumption
offset, and it enters in one place only: $S_{\min}=5\sigma/A$ closes on all
\PbSminClosed{} catalogued rows to \PbSminWorstRes{} in relative terms. On
\PbSminTaut{} rows whose products are not retained $A$ is back-solved from
$5\sigma/S_{\min}$ and flagged as such, so there the closure is arithmetic;
the test is the other \PbSminIndep{} rows, where $A$ comes from the search's
own stored geometry and closes to \PbSminIndepRes. Nothing else needs
correcting: the statistic of
Eq.~\ref{eq:tstar}, the control ranks and the attribution are \emph{ratios}
formed inside one window, in which $A$ divides out identically.

\emph{The star's apparent direction includes the annual parallax, and the
amplitude a compact source retains when that is left out of the assumed phase
centre is measured rather than modelled.} A small fraction of a synthesised
beam attenuates a compact source and does nothing else, while a displacement
approaching a beam puts the phase wrong on the longest baselines, so the size
of the displacement decides whether a scalar correction is legitimate at all.
Re-extracting the \SensNReextWin{} windows of \SensNReextBlock{} blocks
furthest out at the full apparent position -- proper motion, parallax and
radial velocity where catalogued -- leaves each window's noise unchanged to
\SensReextNoiseMaxPct{} per cent and alters no disposition. It cannot test the
amplitude factor, which describes a source those windows do not hold: the
statistic merely scatters from \SensReTstarLo{} to \SensReTstarHi{} about
\SensReTstarMed{} rather than approaching the \SensPxLargePred{} a source
would give. A carrier injected at the apparent position and recovered from the
\emph{same} visibilities at both does return the normalised dirty beam, on
that block's own baselines and with no model in it. It measures
\SensPxSmallMeas{} against a predicted \SensPxSmallPred{} at
\SensPxSmallBeams{} of a beam. At \SensPxLargeBeams{} of a beam it measures
\SensPxLargeMeas{} against \SensPxLargePred. It
exceeds the model at the larger displacement because maximising over a plane
of trajectories lets a partly decohered carrier still find a cell, so a limit
corrected by the model is shallower than it need be.

\emph{A component of a bound pair does not travel in a straight line, and the
propagation assumes it does.} A Gaia~DR3 five-parameter solution carries a
star's position and its velocity at the catalogue epoch and nothing further,
so a linear propagation omits the curvature of the orbit: zero at that epoch,
growing as the square of the interval, and reaching years across an archive
spanning \OmEpochLo--\OmEpochHi. Of the \OmNStar{} searched positions --
the \OmNCensusStar{} census stars and \OmXStar, whose re-extracted windows
the released per-window catalogue does not carry --
\OmNComponent{} have a companion located inside \OmRadius{} arcsec by
Gaia~DR3 or measured in the double-star catalogues
\citep{wds,orb6}, and \OmNIndOnly{} more
have a disturbed astrometric solution with no companion found. Where the pair
has a published visual orbit the displacement follows from it; where it has
not, the projected separation is a lower bound on the true one and so puts a
ceiling on the orbital acceleration and on the displacement, with the
companion's mass taken from the main-sequence sequence of
\citet{pecaut2013}; and the
Hipparcos--Gaia proper-motion anomaly, which is the departure from uniform
motion measured directly over \OmHGInterval{} years
\citep{kervella2022}, covers \OmNPma{} of
them. The largest displacement anywhere in the survey is \OmWorstArcsec{}
arcsec at \OmWorstStar, \OmWorstBeams{} of that block's synthesised beam,
against the \SensPxLargeBeams{} of a beam at which the response measured
above falls to \SensPxLargePred. The largest for which an orbit rather than a
bound is available is
\OmMeasArcsec{} arcsec at \OmMeasStar, \OmMeasBeams{} of a beam. The deepest
limit of any system moves by at most \OmSysMovedMaxPct{} per cent, those of
\OmSysMovedStars, and the median system limit does not move. A displacement
can only attenuate: it can leave a limit optimistic and it can hide a
carrier, but it cannot manufacture a crossing, and no crossing in the survey
lies more than \OmCrossWorstBeams{} of a beam from its predicted position.
The displacement exceeds a tenth of the synthesised beam for \OmNTenth{}
components in all -- \OmTenthStars{} -- and for \OmFlagStars{} it cannot be
computed at all, the solution being disturbed with no companion located and no
Hipparcos epoch to measure one from; those \OmNFlag{} stars carry
\OmNWinFlag{} of the survey's windows, and their limits are optimistic by an
unmeasured amount.

%% Carried forward: the UV Ceti worked example (flux -> EIRP) is out of this
%% appendix because two of its four numbers had no macro behind them.  If a
%% worked example is wanted back it needs three: the flux density, the implied
%% EIRP and the distance used.  It is the clearest single verification aid the
%% paper could carry.  \label{app:repro} and the transmitter-power benchmark
%% are both resolved: nothing cites the label any more, and exactly one
%% benchmark survives, in Sec. 6.
%% tab:blockledger lives here rather than in Sec. 3: it is a full-width
%% table* worth 0.53 pp of the main text, and Sec. 3 carries its headline
%% rows in prose and cites it.
\begin{table*}
\centering
\scriptsize
\caption{\textbf{The survey ledger, part one: execution blocks and windows.}
Each row carries its own number, and a group sums to the row whose number
heads it. The
primary census is the survey: unless a statement names one of the other sets,
it is computed over that row; rows~3.2 to 3.5 are the archival gap, and no
result is formed on them. Hours are window-hours: on-source time summed
over spectral windows, so the classes partition the hours as they partition
the windows. The telescope on-source time behind the primary census is
\StTelHours\,h, of which \StTelHoursFine\,h carries at least one Class~A
window and \StTelHoursCoarse\,h at least one Class~B window -- figures that
overlap and must not be added. Confirmed planet counts here and in the
target list come from the \SfHostCatalogue{} \SfHostTable{} table
\citep{NEA2013}, accessed \SfHostDate. Part two, Table~\ref{tab:counts},
carries the events and the confirmation extents.}
\label{tab:blockledger}
%% GENERATED by blockledger_v411.py -- do not hand-edit.
\begin{tabular}{@{}r@{~}p{0.33\textwidth}rrrrr@{}}
\hline
 & Set of execution blocks & Blocks & Windows & Stars & Systems & Hours \\
\hline
1 & Archival scope: blocks the query exposes & 656 & -- & -- & -- & -- \\
1.1 & \hspace{0.8em}\ldots{} processed with the frozen pipeline & 479 & -- & -- & -- & -- \\
1.2 & \hspace{0.8em}\ldots{} not processed at the freeze & 177 & -- & -- & -- & -- \\
\hline
2 & Processed with the frozen pipeline$^{c}$ & 484 & -- & -- & -- & -- \\
2.1 & \hspace{0.8em}\textbf{Primary census} & 407 & 1663 & 90 & 83 & 1098 \\
2.1.1 & \hspace{1.6em}Class~A, fine channels$^{a}$ & -- & 405 & 66 & 61 & 228 \\
2.1.2 & \hspace{1.6em}Class~B, coarse channels$^{a}$ & -- & 1258 & 89 & 82 & 870 \\
2.2 & \hspace{0.8em}Reserved hold-out & 77 & 315 & 36 & -- & 214 \\
\hline
3 & Not processed at the freeze: what became of each & 177 & -- & -- & -- & -- \\
3.1 & \hspace{0.8em}Calibrated and searched & 144 & -- & -- & -- & -- \\
3.1.1 & \hspace{1.6em}\textbf{Archival extension}, reported & 115 & 463 & 16 & 16 & 269 \\
3.1.1.1 & \hspace{2.4em}Class~A, fine channels$^{a}$ & -- & 67 & 11 & -- & 48 \\
3.1.2 & \hspace{1.6em}Searched after the extension was frozen & 8 & 32 & -- & -- & -- \\
3.1.3 & \hspace{1.6em}$\epsilon$\,Eri, both pointings$^{d}$ & 8 & 32 & -- & -- & -- \\
3.1.4 & \hspace{1.6em}No window reached the export & 13 & 0 & -- & -- & -- \\
3.2 & \hspace{0.8em}No calibration in the ALMA delivery & 18 & -- & -- & -- & -- \\
3.3 & \hspace{0.8em}Calibration failed & 11 & -- & -- & -- & -- \\
3.4 & \hspace{0.8em}Not attempted & 3 & -- & -- & -- & -- \\
3.5 & \hspace{0.8em}Working disk insufficient & 1 & -- & -- & -- & -- \\
\hline
4 & Beyond 40\,pc: out-of-sample control & 47 & 198 & 22 & -- & 88 \\
5 & Rank-calibration sample$^{b}$ & 326 & 1322 & 56 & -- & -- \\
\hline
\end{tabular}

{\footnotesize\raggedright $^{a}$The two classes partition the census windows, not its blocks: 255 blocks carry a Class~A window and 392 a Class~B one, and most carry both, so the hours of the two rows partition the windows and not the observing time.\\ $^{b}$Not an out-of-sample set and not an addition to this ledger: every one of its blocks is counted in a row above, in the primary census or in the hold-out, and it is listed apart only to keep it from being counted twice.\\ $^{c}$Panels 1 and 2 overlap rather than add: 479 of these 484 blocks are row 1.1, and the other 5 carry no member observing unit in the archive metadata, so the scope query does not expose them; 4 of the 5 are in the primary census.\\ $^{d}$Every one is an $\epsilon$\,Eri Band~6 block. The archival pointings put the star either at the centre of the field, at a response of 1.00, or well off it, where it is 0.23--0.28 -- inside the 0.22 floor of \S\ref{sec:sample} in every case, which is where the documented Gaussian and the blocked-aperture form part company and is not a retention limit. All of them are searched, each from its own recorded phase centre, and their conventions are in Appendix~\ref{app:conventions}.\\ The same band's 12 windows that fall in the primary census are counted in row~2.1 and carry every census statistic; the on-axis pointings reach a response of 1.00.}

\end{table*}

\begin{table}
\centering
\scriptsize
\caption{\textbf{The survey ledger, part two: events, dispositions and
confirmation.} Every headline count in this paper is read from here, and
\textbf{the itemisation accompanies this paper
as data rather than as part of the article}, machine-readable and carrying all
\MtCountRows{} rows, which is why the numbering here is not consecutive. Three
counts are \MtCrossA{} over three different sets -- the Class~A crossings of
the adopted ledger, those of the per-window extraction catalogue
(\S\ref{sec:sample}), and the
quality-passing crossings. The first two hold different rows, because on some
windows the two read the statistic over different fields -- the whole primary
beam in the ledger and a narrower one in the catalogue
(Appendix~\ref{app:falsealarm}) -- which puts crossings on opposite sides of
the trigger in both directions; the last two hold the same rows,
because the one block failing the quality criterion is also the survey's only
block with a coarse-channel crossing. Row 3.1.1 is what this
paper means by confirmation and the only such count it quotes: it applies the
criterion of \S\ref{sec:summary}, a different observing session and more than
a day, at a frequency the first block covered. Row 2.2.1 is
decomposed system by system in Appendix~\ref{app:falsealarm}.}
\label{tab:counts}
%% GENERATED by counts_r15.py -- do not hand-edit.
%% The accounting rows of the master count table; the itemisation is in
%% the machine-readable file this generator writes beside it.
\begin{tabular}{@{}r@{~}p{0.405\columnwidth}r@{\,\,}p{0.325\columnwidth}@{}}
\hline
 & Count & Value & What it counts \\
\hline
1 & Threshold crossings at $T_\star\ge5$ & \EvNCrossRaw & census and archival tail \\
1.1 & Class~A, fine channels & \MtCrossA & the carrier experiment \\
1.2 & Class~B, coarse channels & \LdgNCrossB & all in the one block failing \S\ref{sec:bdblock} \\
\hline
2 & Quality-passing crossings & \EvNCrossQp & the same rows as 1.1 \\
2.1 & \ldots{} attributed to a transition & \EvNAttrQp & of \MkNCoinc{} coincidences inside the mask in the star's own frame \\
2.2 & \ldots{} unattributed & \EvNUnattrQp & rows 2.2.1 and 2.2.2 \\
2.2.1 & \ldots{} in the primary census & \CeObs & against \CeExp{} expected by chance at the \CeNWin{} windows' own control positions -- a count chance delivers with probability \CeProbNull \\
2.2.2 & \ldots{} in archival-tail blocks & \CeNTailCross & outside that comparison \\
\hline
3 & Class~A systems & \NSysClassA & of \NSystems{} surveyed \\
3.1.1 & Class~A systems with a later block at a frequency the first covered, in another session & \textbf{\ZcConfSysDay} & \textbf{the confirmable set, and the only one quoted}; \ZcConfSysYr{} of them over a year later. Over the census together with the repeat blocks the recurrence test reads \\
3.3 & \ldots{} with no separated second block at a covered frequency & \ZcNoRepSys & confirmation impossible over the \ZcNRecEB{} blocks the recurrence test reads; \RcNoRepSysAll{} of all \NSystems \\
\hline
4 & Class~A union bandwidth & \CvUnionA\,GHz & \DomAIslands{} disjoint islands \\
4.1 & \ldots{} removed by the line mask & \MkMaskAPct\,\% & \MkMaskAGHz\,GHz: the figure the text quotes, against \FrMaskPct{} per cent of the \DnuAB\,GHz archival union \\
\hline
\end{tabular}

\end{table}

%% ★★ 2026-10-08: THE ALMA PROJECT CODES, RESTORED AS A TABLE.
%% They were taken out of the Acknowledgements as padding and then lived "in
%% the data release"; there is no data release, so they were nowhere.  ALMA's
%% data-use policy expects the codes to be traceable from the publication, so
%% this is not optional -- but an appendix table satisfies the policy without
%% putting \NProjCodes{} codes back into the front matter, which is what was
%% objected to.  The fragment is generated by projcodes_v400.py and is the
%% same list that resolves every block in Table~\ref{tab:blockledger}.
\begin{table}
\centering
\scriptsize
\caption{\textbf{The \NProjCodes{} ALMA proposals the searched execution
blocks belong to}, in the standard
\mbox{ADS/JAO.ALMA\#\textit{code}} form. Every block of the primary census
belongs to one of these, and the blocks themselves are retrieved from the
ALMA Science Archive by the execution-block identifiers
the accompanying ledger file gives for every crossing.\ProjCodeUnresClause}
\label{tab:projcodes}
%% GENERATED by projcodes_v400.py -- do not hand-edit.
\mbox{2012.1.00142.S}, \mbox{2012.1.00198.S}, \mbox{2012.1.00268.S}, \mbox{2012.1.00385.S}, \mbox{2012.1.00437.S}, \mbox{2012.1.00993.S}, \mbox{2013.1.00359.S}, \mbox{2013.1.00588.S}, \mbox{2013.1.00645.S}, \mbox{2013.1.01147.S}, \mbox{2013.1.01214.S}, \mbox{2015.1.00307.S}, \mbox{2015.1.00633.S}, \mbox{2015.1.00783.S}, \mbox{2015.1.01015.S}, \mbox{2015.1.01260.S}, \mbox{2016.1.00104.S}, \mbox{2016.1.00358.S}, \mbox{2016.1.00515.S}, \mbox{2016.1.00637.S}, \mbox{2016.1.00803.S}, \mbox{2016.1.00817.S}, \mbox{2016.1.00878.S}, \mbox{2016.1.00880.S}, \mbox{2016.2.00015.S}, \mbox{2016.2.00200.S}, \mbox{2016.A.00013.S}, \mbox{2017.1.00167.S}, \mbox{2017.1.00200.S}, \mbox{2017.1.00215.S}, \mbox{2017.1.00561.S}, \mbox{2017.1.00698.S}, \mbox{2017.1.00704.S}, \mbox{2017.1.00786.S}, \mbox{2017.1.00986.S}, \mbox{2017.1.01583.S}, \mbox{2017.1.01644.S}, \mbox{2018.1.00045.S}, \mbox{2018.1.00072.S}, \mbox{2018.1.00126.S}, \mbox{2018.1.00461.S}, \mbox{2018.1.00470.S}, \mbox{2018.1.00947.S}, \mbox{2018.1.01149.S}, \mbox{2018.1.01545.S}, \mbox{2018.1.01833.S}, \mbox{2019.1.00189.S}, \mbox{2019.1.00696.S}, \mbox{2019.1.01443.T}, \mbox{2019.1.01517.S}, \mbox{2019.1.01603.S}, \mbox{2019.1.01681.S}, \mbox{2019.2.00208.S}, \mbox{2021.1.00485.S}, \mbox{2021.1.01209.S}, \mbox{2021.2.00097.S}, \mbox{2022.1.00134.S}, \mbox{2022.1.00338.L}, \mbox{2022.1.00793.S}, \mbox{2022.1.01163.S}, \mbox{2023.1.00390.S}, \mbox{2023.1.00487.S}, \mbox{2024.1.00165.S}, \mbox{2024.1.00595.S}, \mbox{2024.1.00722.S}, \mbox{2024.1.01095.S}, \mbox{2024.A.00040.S}

\end{table}
